\documentclass[../main]{subfiles}
\usepackage{tikz}
\usetikzlibrary{arrows.meta,decorations.markings,calc,patterns}
\begin{document}

\chapter{Uso de Arduino con Servos y Motores Paso a Paso}

\section{Control de Servos y Motores Paso a Paso con Arduino y el Shield DK Electronics V1.0}

\subsection{Introducción a Arduino y el Shield DK Electronics V1.0}

Arduino es una plataforma de prototipado electrónico de código abierto que permite crear proyectos interactivos. El shield DK Electronics V1.0 de Adafruit\footnote{Adafruit: Empresa especializada en la creación de hardware y software de código abierto para proyectos de electrónica y computación física.} es un accesorio que facilita la conexión y control de servos y motores paso a paso.

\info{Descripción del Shield DK Electronics V1.0}{
El shield DK Electronics V1.0 está diseñado para ser montado directamente sobre un Arduino y permite controlar hasta dos motores paso a paso\footnote{Motor paso a paso: Motor que se mueve en pasos discretos, permitiendo un control preciso de su posición.} o cuatro servos\footnote{Servo: Dispositivo que utiliza un motor para posicionar su eje en un ángulo específico basado en una señal de control.}. Este shield incluye un controlador de motores L293D que maneja la potencia de los motores.
}

\subsection{Control de un Servo}

Para controlar un servo con Arduino, se utiliza la biblioteca `Servo.h`, que permite enviar señales PWM\footnote{PWM (Pulse Width Modulation): Técnica de modulación de ancho de pulso utilizada para controlar dispositivos electrónicos.} al servo para posicionarlo en el ángulo deseado.

\begin{equation}
\theta = \dfrac{\text{duty cycle} \times 180}{100}
\end{equation}

Siendo:
\begin{itemize}
    \item $\theta$: Ángulo del servo en grados.
    \item \text{duty cycle}: Ciclo de trabajo de la señal PWM en \%.
\end{itemize}

\ejemplo{Ejemplo: Control de un Servo}{Supongamos que queremos posicionar un servo a 90 grados. Utilizaremos la función `write` de la biblioteca `Servo`:

\begin{verbatim}
#include <Servo.h>

Servo miServo;

void setup() {
  miServo.attach(9); // Conecta el servo al pin 9
}

void loop() {
  miServo.write(90); // Posiciona el servo a 90 grados
  delay(1000); // Espera 1 segundo
}
\end{verbatim}
}

\subsection{Control de un Motor Paso a Paso}

Para controlar un motor paso a paso con el shield DK Electronics V1.0, se utiliza la biblioteca `Stepper.h`. Este shield permite controlar los motores de manera precisa, especificando el número de pasos a moverse.

\begin{equation}
N = \dfrac{\text{Grados a mover} \times \text{Pasos por revolución}}{360}
\end{equation}

Siendo:
\begin{itemize}
    \item $N$: Número de pasos.
    \item \text{Pasos por revolución}: Número de pasos que el motor necesita para completar una revolución completa.
\end{itemize}

\ejemplo{Ejemplo: Control de un Motor Paso a Paso}{Supongamos que queremos mover un motor paso a paso 90 grados, y que el motor tiene 200 pasos por revolución:

\begin{verbatim}
#include <Stepper.h>

const int pasosPorRevolucion = 200;
Stepper miMotor(pasosPorRevolucion, 8, 9, 10, 11);

void setup() {
  miMotor.setSpeed(60); // Establece la velocidad a 60 RPM
}

void loop() {
  miMotor.step(50); // Mueve el motor 50 pasos (90 grados)
  delay(1000); // Espera 1 segundo
}
\end{verbatim}
}

\section{Cuestionario}

\actividad{Responder el siguiente cuestionario}
\begin{enumerate}
    \item ¿Qué es Arduino y para qué se utiliza?
    \item ¿Qué es un shield en el contexto de Arduino?
    \item Explica el funcionamiento de un servo.
    \item ¿Cómo se controla un servo con Arduino?
    \item Define qué es un motor paso a paso y sus aplicaciones.
    \item ¿Qué ventajas ofrece el uso del shield DK Electronics V1.0?
    \item ¿Qué es una señal PWM y cómo se utiliza en el control de servos?
    \item ¿Cómo se calcula el número de pasos para mover un motor paso a paso?
    \item ¿Qué es el ciclo de trabajo en una señal PWM?
    \item ¿Qué función tiene el controlador de motores L293D en el shield DK Electronics?
\end{enumerate}

\section{Problemas}

\actividad{Resolver los siguientes problemas}
\begin{enumerate}
    \item Calcular el número de pasos necesarios para mover un motor paso a paso 45 grados, si el motor tiene 200 pasos por revolución.
    \item Determinar el ciclo de trabajo necesario para posicionar un servo a 135 grados.
    \item Si un motor paso a paso tiene 400 pasos por revolución, ¿cuántos pasos se necesitan para moverlo 180 grados?
    \item Controlar un servo a 60 grados utilizando la biblioteca `Servo.h`. ¿Cuál es el código necesario?
    \item Si el ciclo de trabajo de una señal PWM es 50\%, ¿a qué ángulo se posicionará el servo?
    \item Calcular el número de pasos para mover un motor paso a paso 270 grados, con 200 pasos por revolución.
    \item ¿Cuál es la velocidad en RPM necesaria para mover un motor paso a paso 100 pasos en 2 segundos?
    \item Determinar el número de pasos para mover un motor paso a paso 360 grados, con 400 pasos por revolución.
    \item Si un servo se posiciona a 90 grados con un ciclo de trabajo del 50\%, ¿qué ciclo de trabajo se necesita para 45 grados?
    \item Escribe un programa en Arduino para mover un motor paso a paso 180 grados a una velocidad de 30 RPM.
\end{enumerate}

\end{document}
